% Lab 5 — Additional Constraints（原书 PDF p.71–82）
\bichapter{Lab 5：额外约束}{Lab 5: Additional Constraints}

\begin{objectiveBox}
完成本实验后，你应能够：
\begin{itemize}
  \item 应用用户指定的 annotated delay 探索 latch 的 time borrowing（时间借用）；
  \item 调试时钟路径上 non-unate 单元相关的 \textbf{PTE-070} 信息消息。
\end{itemize}
\end{objectiveBox}

\begin{durationBox}
约 \textbf{30} 分钟。
\end{durationBox}

\section{概述}
\label{lab5:overview}

\noindent\textit{Overview}

本实验涵盖两类额外约束场景：调试时钟网络上的 non-unate 路径警告，以及借助 annotated delay 探索 latch time borrowing。

\subsection{相关文件与目录}
\label{lab5:files}

\noindent\textit{Relevant Files and Directories}

本实验全部文件位于主目录下的 \cmd{lab5\_additional} 目录中。

\begin{tabularx}{\textwidth}{@{}l X@{}}
\toprule
\textbf{路径 / 文件} & \textbf{说明} \\
\midrule
\cmd{lab5\_additional/} & 当前工作目录 \\
\cmd{orca\_savesession/} & ORCA 保存会话 \\
\cmd{RUN.tcl} & 创建 \cmd{orca\_savesession} 的脚本 \\
\cmd{logs/} & 运行脚本产生的日志文件 \\
\cmd{.synopsys\_pt.setup} & PT setup 文件 \\
\bottomrule
\end{tabularx}

\subsection{答案 / 解答}
\label{lab5:answers-note}

\noindent\textit{Answers / Solutions}

本实验手册在每个 Lab 后半部分给出全部问题的答案与解答。
若答题需要帮助，或希望核对结果，请查阅本 Lab 后部的解答章节。

\section{操作说明}
\label{lab5:instructions}

\noindent\textit{Instructions}

\subsection{任务 1：调试 PTE-070 信息消息}
\label{lab5:task1}

\noindent\textit{Task 1. Debug PTE-070 Information Messages}

\begin{enumerate}
  \item 从 \cmd{lab5\_additional} Unix 目录启动 PrimeTime。
        恢复保存在 \cmd{./orca\_savesession} 中的会话。

  \item 下列为时钟网络上 non-unate 路径的完整消息。
        下一步将要求你生成经过该 pin 的时序报告。
        为便于复制粘贴、避免拼写错误，可在另一终端窗口的日志文件中查找该消息，
        或在 PrimeTime 内用 Unix 命令 \cmd{grep}：
\begin{lstlisting}
# 来自 ./logs/run.log
Information: A non-unate path in clock network detected.
Propagating both inverting and noninverting senses of clock
"SDRAM_CLK" from pin
'I_ORCA_TOP/I_SDRAM_IF/sd_mux_dq_out_0/Z'. (PTE-070)
pt_shell> sh grep -A 1 -B 1 PTE-070 logs/run.log
\end{lstlisting}

\begin{noteBox}
命令 \cmd{sh}（或 \cmd{exec}）允许在 PrimeTime shell 内执行 Unix 命令。
\end{noteBox}

  \item 为 setup 生成经过上述 pin 的时序报告，并回答下列问题。
        下列别名已在 PrimeTime setup 文件中创建，
        将借助 SolvNet Doc Id 014947 上的 \cmd{view} 工具在带滚动条的弹出窗口中生成时序报告：
\begin{lstlisting}
pt_shell> vrt -through <through pin>
\end{lstlisting}

\begin{questionBox}
\textbf{问题 1.}\ 你用时序报告中哪些行来确认这是 setup 报告，且路径起点为时钟 \cmd{SDRAM\_CLK} 的 source？\\
\textbf{问题 2.}\ 该时序报告如何确认警告中的 pin 位于 data path（即时钟 source 被当作数据路径使用并约束），而非 clock path？\\
\textbf{问题 3.}\ 经过上述 pin 传播的是哪种 sense（positive unate 还是 negative unate）？
在时序报告中查找标示该 pin 的小箭头。
\end{questionBox}

  \item 再生成至少一份时序报告，展示经过该 pin 的 negative unate 时序弧。

\begin{questionBox}
\textbf{问题 4.}\ 你用时序报告中哪些行来确认这是 setup 报告、路径起点为 \cmd{SDRAM\_CLK} 的 source，
且该 pin 的时序弧为 negative unate？\\
\textbf{问题 5.}\ 说明为何对这些时序路径可忽略（并抑制）该警告。
\end{questionBox}

  \item 不要退出 PrimeTime。
\end{enumerate}

\subsection{任务 2：探索 Time Borrow 与 Latch}
\label{lab5:task2}

\noindent\textit{Task 2. Explore Time Borrow and Latches}

本设计中只有一个 latch。

\begin{enumerate}
  \item 用下列命令找到它（可键入少量字符后按 \textbf{Tab} 键补全）：
\begin{lstlisting}
pt_shell> all_registers -level_sensitive
pt_shell> !! -clock_pins
pt_shell> all_registers -level_sensitive -data_pins
\end{lstlisting}

\begin{questionBox}
\textbf{问题 6.}\ 该 latch 的 clock pin 名称是什么？\\
\textbf{问题 7.}\ 三个 data pin 的名称分别是什么？
\end{questionBox}

  \item 从该 latch 生成 setup 时序报告（须明确以 clock pin 为起点，而非仅用单元名）。
        本实验将称该报告为「path segment \#2」。

\begin{noteBox}
该 latch 在 ORCA 设计中的功能是产生时钟门控信号，用于打开/关闭时钟 \cmd{SYS\_CLK}。
\end{noteBox}

\begin{questionBox}
\textbf{问题 8.}\ 你如何知道该 latch 未从上一级经历 time borrow？
\end{questionBox}

  \item 为上一级生成时序报告（本实验称为「path segment \#1」）。
        以 latch 的 D 输入 pin 作为该时序路径的终点：
\begin{lstlisting}
pt_shell> report_timing -to I_ORCA_TOP/I_BLENDER/latched_clk_en_reg/D
\end{lstlisting}

\begin{questionBox}
\textbf{问题 9.}\ path segment \#1 还能再花多少时间才会开始从 path segment \#2 借用时间？
\end{questionBox}

  \item 通过标注 4\,ns 的 net delay，强制 path segment \#1 从 path segment \#2 借用时间：
\begin{lstlisting}
# 复制粘贴以避免 pin 名拼写错误
pt_shell> set_annotated_delay -net 4 \
    -to I_ORCA_TOP/I_BLENDER/latched_clk_en_reg/D
\end{lstlisting}

  \item 再次生成 path segment \#1 的时序报告（可用上下箭头在历史事件列表中滚动）。

\begin{questionBox}
\textbf{问题 10.}\ path segment \#1 从 path segment \#2 借用了多少时间？\\
\textbf{问题 11.}\ path segment \#1 的 slack 是多少？
\end{questionBox}

  \item 重新生成 path segment \#2 的时序报告。在此之前须执行完整 timing update：
\begin{lstlisting}
pt_shell> update_timing -full
\end{lstlisting}

\begin{noteBox}
此时路径起点将为 latch 的 D pin（而非先前使用的 clock pin），
因为你关注的是包含 time borrow 的时序路径。
\end{noteBox}

\begin{lstlisting}
pt_shell> report_timing -from \
    I_ORCA_TOP/I_BLENDER/latched_clk_en_reg/D
\end{lstlisting}

\begin{questionBox}
\textbf{问题 12.}\ 给予 path segment \#1 的时间是否符合你的预期？
\end{questionBox}

  \item 修改 latch 行为以实现 transparency（透明）：
        即当数据在时钟 opening 与 closing 边之间到达时，latch 保持透明。
\begin{lstlisting}
set_app_var timing_enable_through_paths true
\end{lstlisting}

  \item 重复到 latch D pin 的时序报告。注意即使 latch 处于透明状态，仍可将 D pin 指定为终点：
\begin{lstlisting}
report_timing -to I_ORCA_TOP/I_BLENDER/latched_clk_en_reg/D
\end{lstlisting}

\begin{questionBox}
\textbf{问题 13.}\ 起点（startpoint）是什么？
\end{questionBox}

  \item 从刚识别的起点做时序报告：
\begin{lstlisting}
report_timing -from I_ORCA_TOP/I_PARSER/blender_clk_en_reg/CP
\end{lstlisting}

\begin{questionBox}
\textbf{问题 14.}\ 路径终点是什么？\\
\textbf{问题 15.}\ transparency 的 open 与 close 边分别是多少？\\
\textbf{问题 16.}\ 数据是否在两者之间到达？\\
\textbf{问题 17.}\ 是否发生了 time borrowing？\\
\textbf{问题 18.}\ slack 是否为正？
\end{questionBox}

  \item 退出 PrimeTime。
\end{enumerate}

\noindent 至此 Lab~5 完成。请返回课堂讲义。

\section{答案 / 解答}
\label{lab5:solutions}

\noindent\textit{Answers / Solutions}

\begin{answerBox}
\textbf{问题 1.}
\begin{lstlisting}
pt_shell> report_timing -through \
    I_ORCA_TOP/I_SDRAM_IF/sd_mux_dq_out_0/Z
Startpoint: sdr_clk (clock source 'SDRAM_CLK')
Endpoint:   sd_DQ[0] (output port clocked by SD_DDR_CLK)
Path Group: SD_DDR_CLK
Path Type:  max
\end{lstlisting}
\texttt{Path Type: max} 表示 setup；\texttt{Startpoint: sdr\_clk (clock source 'SDRAM\_CLK')}
表明起点为 \cmd{SDRAM\_CLK} 的 source。
\end{answerBox}

\begin{answerBox}
\textbf{问题 2.}\ 若未生成报告（``path is unconstrained''），则该 pin 位于 clock path。
由于生成了时序报告，该 pin 位于 \textbf{data path}。
\end{answerBox}

\begin{answerBox}
\textbf{问题 3.}\ 经过该 pin 的是 \textbf{positive unate} 时序弧（fall to fall）。
数据通路中的相关行：
\begin{lstlisting}
I_ORCA_TOP/I_SDRAM_IF/bufd7G5B2136/Z (bufd7)  0.1634 & 4.8748 f
I_ORCA_TOP/I_SDRAM_IF/sd_mux_dq_out_0/Z (mx02d4) <- 0.5431 & 5.4179 f
\end{lstlisting}
箭头 \texttt{<-} 表示 negative unate；无箭头表示 positive unate。
\end{answerBox}

\begin{answerBox}
\textbf{问题 4.}
\begin{lstlisting}
pt_shell> report_timing -rise_through \
    I_ORCA_TOP/I_SDRAM_IF/sd_mux_dq_out_0/Z
Startpoint: sdr_clk (clock source 'SDRAM_CLK')
Endpoint:   sd_DQ[0] (output port clocked by SD_DDR_CLK)
Path Group: SD_DDR_CLK
Path Type:  max
I_ORCA_TOP/I_SDRAM_IF/bufd7G5B2136/Z (bufd7)  0.1634 & 4.8748 f
I_ORCA_TOP/I_SDRAM_IF/sd_mux_dq_out_0/Z (mx02d4) <- 0.4928 & 5.3676 r
\end{lstlisting}
\texttt{Path Type: max} 与起点确认 setup 及 \cmd{SDRAM\_CLK} source；
箭头 \texttt{<-} 表示该 pin 为 \textbf{negative unate} 弧。
\end{answerBox}

\begin{answerBox}
\textbf{问题 5.}\ 该消息表示传播时钟经过该 MUX 时将同时使用反相与非反相 sense——这是默认行为。
但由于该时钟被当作数据使用，PrimeTime 实际上会传播两种 sense（positive 与 negative unate），
即使在较旧版本的 PrimeTime 中也是如此。这正是期望的行为，因此该信息消息可忽略并抑制。
\end{answerBox}

\begin{answerBox}
\textbf{问题 6.}
\begin{lstlisting}
pt_shell> all_registers -level_sensitive -clock_pins
{"I_ORCA_TOP/I_BLENDER/latched_clk_en_reg/EN"}
\end{lstlisting}
\end{answerBox}

\begin{answerBox}
\textbf{问题 7.}
\begin{lstlisting}
pt_shell> all_registers -level_sensitive -data_pins
{"I_ORCA_TOP/I_BLENDER/latched_clk_en_reg/D",
 "I_ORCA_TOP/I_BLENDER/latched_clk_en_reg/SC",
 "I_ORCA_TOP/I_BLENDER/latched_clk_en_reg/SD"}
\end{lstlisting}
\end{answerBox}

\begin{answerBox}
\textbf{问题 8.}\ 若该 latch 正在经历 time borrow，报告中会出现一行说明给予起点（即上一级）的时间。
本报告中\textbf{没有}该行，故未经历 time borrow。

生成 path segment \#2 的命令：
\begin{lstlisting}
pt_shell> report_timing -from \
    I_ORCA_TOP/I_BLENDER/latched_clk_en_reg/EN
\end{lstlisting}
\end{answerBox}

\begin{answerBox}
\textbf{问题 9.}\ 还可再花 \textbf{3.465\,ns} 才会开始从 path segment \#2 借用时间
（等于正 slack）。
\end{answerBox}

\begin{answerBox}
\textbf{问题 10.}\ path segment \#1 借用了 \textbf{0.5306\,ns}（在时序报告的 data required time 段中注明）。
\end{answerBox}

\begin{answerBox}
\textbf{问题 11.}\ slack 为 \textbf{0}。
PrimeTime 恰好借用所需时间，使 slack 等于零。
\end{answerBox}

\begin{answerBox}
\textbf{问题 12.}\ 是——path segment \#2 时序报告中给予起点的时间，
与 path segment \#1 时序报告中借用的 time borrow 一致。
\begin{lstlisting}
pt_shell> report_timing -from I_ORCA_TOP/I_BLENDER/latched_clk_en_reg/D
pt_shell> report_timing -to I_ORCA_TOP/I_BLENDER/latched_clk_en_reg/D
\end{lstlisting}
\end{answerBox}

\begin{answerBox}
\textbf{问题 13.}\ \textbf{I\_ORCA\_TOP/I\_PARSER/blender\_clk\_en\_reg/CP}。
\end{answerBox}

\begin{answerBox}
\textbf{问题 14.}\ \textbf{I\_ORCA\_TOP/I\_BLENDER/U794/A}（gating check 终点）。
\end{answerBox}

\begin{answerBox}
\textbf{问题 15.}\ transparency open 与 close 边分别为 \textbf{6.6454} 与 \textbf{10.4433}。
\end{answerBox}

\begin{answerBox}
\textbf{问题 16.}\ 是——数据在 \textbf{7.1759} 到达，介于 open 与 close 边之间。
\end{answerBox}

\begin{answerBox}
\textbf{问题 17.}\ \textbf{否}，未发生 time borrowing。
\end{answerBox}

\begin{answerBox}
\textbf{问题 18.}\ \textbf{是}，slack 为正（\textbf{2.1398\,ns}）。

完整透明 latch 时序报告：
\begin{lstlisting}
pt_shell> report_timing -from \
    I_ORCA_TOP/I_PARSER/blender_clk_en_reg/CP
Startpoint: I_ORCA_TOP/I_PARSER/blender_clk_en_reg
  (rising edge-triggered flip-flop clocked by SYS_CLK)
Endpoint: I_ORCA_TOP/I_BLENDER/U794
  (rising clock gating-check end-point clocked by SYS_CLK)
Path Group: **clock_gating_default**
Path Type: max
...
transparency open edge   6.6454
transparency close edge 10.4433
available borrow at through pin 3.2674
...
slack (MET) 2.1398
\end{lstlisting}
\end{answerBox}
